\section{Actual finite chart relations have homogeneous polynomial lifts}

Fix a genuine point-chart pivot $e$ and direction parameters $b$.
Write $\delta_{e,b}$ for the actual polynomial homomorphism sending
$x_e$ to $1$ and $x_i$ to $b_i+x_i$ for $i\ne e$. Every ring below
is the literal finite polynomial ring over $\mathbb C$.

\begin{lemma}[Explicit actual rehomogenization]
\label{hir:actual-polynomial-rehomogenization}
\uses{hir:actual-point-polynomial-dehomogenization}
\lean{Hironaka.HironakaActualPolynomialRehomogenization.actualPolynomialRehomogenization,Hironaka.HironakaActualPolynomialRehomogenization.actualPolynomialRehomogenization_homogeneous,Hironaka.HironakaActualPolynomialRehomogenization.actualPointDehomogenization_rehomogenization}
\leanok
Let $q$ be a complex polynomial of total degree at most $m$. Define
\[
 H_{e,b,m}(q)=\sum_{\alpha\in\operatorname{supp}q}
 x_e^{m-|\alpha|}\,q_\alpha
 \prod_i(x_i-b_i x_e)^{\alpha_i}.
\]
This polynomial is homogeneous of degree $m$. If every exponent in
$\operatorname{supp}q$ has $\alpha_e=0$, then
$\delta_{e,b}H_{e,b,m}(q)=q$ as whole finite polynomials.
\end{lemma}
\begin{proof}
Each linear form $x_i-b_i x_e$ is homogeneous of degree one, including
when that form is zero. The monomial product therefore has degree
$|\alpha|$. Every support exponent satisfies $|\alpha|\le\deg q\le m$,
so multiplication by $x_e^{m-|\alpha|}$ makes each summand homogeneous
of degree exactly $m$. Their finite sum has that same degree, with zero
polynomials allowed. Under the actual dehomogenization the pivot power
is one. Each other linear form maps to $x_i$. The pivot linear form
might instead map to the constant $1-b_e$, but its exponent is zero
by the stated support condition. Thus each whole summand maps to
$q_\alpha x^\alpha$. The genuine finite support expansion of $q$
proves the whole polynomial equality.
\end{proof}

\begin{lemma}[Actual dehomogenization is injective in each fixed homogeneous degree]
\label{hir:actual-homogeneous-polynomial-dehomogenization-injective}
\uses{hir:actual-point-blowup-homogeneous-polynomial-factor,hir:actual-formal-blowup-all-directions-whole-power-reflection}
\lean{Hironaka.HironakaActualPolynomialRehomogenization.actualPointDehomogenization_homogeneous_injective}
\leanok
If $P,Q$ are homogeneous polynomials of the same degree $D$ and
$\delta_{e,b}P=\delta_{e,b}Q$, then $P=Q$.
\end{lemma}
\begin{proof}
Include both polynomials in the actual formal series ring. The genuine
whole homogeneous point-blowup factorization gives
$\pi^*P=x_e^D\delta_{e,b}P$ and
$\pi^*Q=x_e^D\delta_{e,b}Q$. Their right sides agree. The actual formal
blowup homomorphism is injective for every genuine chart direction,
so the whole formal inclusions of $P,Q$ agree. The genuine polynomial
inclusion is injective, hence the original whole polynomials agree.
\end{proof}

\begin{theorem}[Every finite pivot-independent chart relation lifts to one homogeneous relation]
\label{hir:actual-finite-polynomial-chart-relation-rehomogenization}
\uses{hir:actual-polynomial-rehomogenization,hir:actual-homogeneous-polynomial-dehomogenization-injective}
\lean{Hironaka.HironakaActualPolynomialRehomogenization.homogeneous_polynomial_relation_lift_at_degree,Hironaka.HironakaActualPolynomialRehomogenization.exists_homogeneous_polynomial_relation_lift}
\leanok
Let $f_i$ be a finite family of actual complex polynomials homogeneous
of degrees $w_i$. Let $q_i$ be actual polynomials independent of the
pivot coordinate, in the literal support sense of the first lemma.
If $\sum_iq_i\delta_{e,b}f_i=0$, there are one integer $D\ge w_i$
and homogeneous actual polynomials $P_i$ of degrees $D-w_i$ such that
\[
 \delta_{e,b}P_i=q_i,\qquad \sum_iP_if_i=0.
\]
No relation lift, source ideal, analytic flatness or convergence
certificate is supplied.
The same conclusion holds at any prescribed $D$ satisfying the natural
finite bounds $w_i+\deg q_i\le D$.
\end{theorem}
\begin{proof}
Choose the actual finite integer $D=\sum_i(w_i+\deg q_i)$.
Every summand is nonnegative, so $D\ge w_i+\deg q_i$ for every index.
Thus $\deg q_i\le D-w_i$, and the first lemma defines
$P_i=H_{e,b,D-w_i}(q_i)$ with the required degree and exact
whole dehomogenization. Each product $P_if_i$ is homogeneous of
degree $(D-w_i)+w_i=D$, so the entire finite sum is homogeneous
of degree $D$. Dehomogenization is a ring homomorphism; its image
is exactly the assumed zero chart relation. The preceding fixed-degree
injectivity lemma, comparing with the zero polynomial of degree $D$,
forces the whole original sum to be zero. The construction also covers
an empty family, when $D=0$ and both sums are zero.
\end{proof}
